\documentclass[../DoAn.tex]{subfiles}
\begin{document}

\section{Đặt vấn đề và tính cấp thiết của đề tài}
\label{section:1.1}
Sự bứt phá của ngành giao thông vận tải và mạng lưới logistics toàn cầu đang định hình lại cấu trúc kinh tế của thế giới hiện đại. Sự gia tăng theo cấp số nhân của các đội xe thương mại, dịch vụ gọi xe và phương tiện cá nhân mang lại lợi ích to lớn về lưu thông hàng hóa và con người. Tuy nhiên, đi kèm với sự phát triển này là những thách thức sinh tử liên quan đến an toàn giao thông đường bộ. Trạng thái tâm sinh lý và mức độ nhận thức của người điều khiển phương tiện (tài xế) đóng vai trò trung tâm trong chuỗi nhân quả dẫn đến các thảm kịch trên cao tốc và đường đô thị. Những sự cố này hiếm khi xuất phát từ lỗi cơ khí thuần túy, mà phần lớn bắt nguồn từ sự suy giảm nhận thức thoáng qua hoặc những phản ứng cảm xúc cực đoan vượt tầm kiểm soát của con người.

Thực tiễn này đòi hỏi một sự chuyển dịch mô hình (paradigm shift) từ các biện pháp an toàn thụ động sang các hệ thống an toàn chủ động có khả năng nhận thức ngữ cảnh. Việc ứng dụng Trí tuệ nhân tạo (AI) và Thị giác máy tính (Computer Vision) vào không gian cabin xe không chỉ là một giải pháp công nghệ đơn thuần mà đã trở thành một yêu cầu mang tính chiến lược, được định hướng bởi quy luật cung cầu của thị trường lẫn các khung pháp lý quốc tế khắt khe.

\subsection{Rủi ro tai nạn do trạng thái tâm sinh lý của tài xế}
Để thấu hiểu tính cấp thiết của công nghệ giám sát, việc mổ xẻ nguyên nhân gốc rễ của tai nạn giao thông dưới góc độ dữ liệu dịch tễ học và tâm lý học hành vi là vô cùng quan trọng. Theo thống kê của Tổ chức Y tế Thế giới (WHO), tai nạn giao thông đường bộ tước đi sinh mạng của xấp xỉ 1,3 triệu người mỗi năm, đồng thời là nguyên nhân tử vong hàng đầu đối với thanh thiếu niên và người trưởng thành trẻ tuổi (từ 5 đến 29 tuổi) trên toàn cầu. Một tỷ lệ lớn các vụ tai nạn nghiêm trọng này có nguyên nhân trực tiếp từ trạng thái tâm sinh lý không ổn định của người điều khiển phương tiện. Các yếu tố rủi ro cốt lõi bao gồm:

\begin{enumerate}
    \item \textbf{Trạng thái buồn ngủ và mệt mỏi thể chất (Fatigue/Drowsiness):}\\
    Việc vận hành phương tiện liên tục trong nhiều giờ, đặc biệt vào ban đêm, rạng sáng, hoặc trên các tuyến đường cao tốc đơn điệu, dễ dẫn đến hiện tượng suy giảm nhận thức cấp tính. Các nghiên cứu đánh giá rủi ro toàn cầu chỉ ra rằng tình trạng buồn ngủ khi lái xe đóng góp từ 3\% đến hơn 30\% tổng số các vụ tai nạn giao thông. Sự xao nhãng này thường liên quan đến tình trạng thiếu ngủ kéo dài, làm gián đoạn nhịp sinh học (circadian misalignment), đặc biệt phổ biến ở các tài xế có xu hướng hoạt động vào buổi sáng (morning chronotype) nhưng buộc phải làm việc vào chiều muộn hoặc đêm khuya. Đáng chú ý, hơn 20\% tài xế thừa nhận từng cảm thấy cần phải dừng xe khẩn cấp ít nhất một lần do không thể cưỡng lại cơn buồn ngủ.
    
    Biểu hiện nguy hiểm nhất của trạng thái này là sự xuất hiện của các cơn ngủ gật ngắn (microsleep). Đây là các chu kỳ mất ý thức tạm thời kéo dài chỉ từ 1 đến 3 giây. Ở vận tốc 100 km/h, một phương tiện trong 2 giây mất ý thức có thể di chuyển vô định xấp xỉ 55 mét. Trong khoảng thời gian mù này, phương tiện có thể đi lệch làn đường hoặc lao thẳng vào chướng ngại vật. Tai nạn do mệt mỏi và buồn ngủ có khả năng gây ra chấn thương nghiêm trọng và tỷ lệ tử vong cao hơn hẳn các loại tai nạn khác, do người lái xe không có bất kỳ hành động phanh hay đánh lái phòng tránh nào trước khoảnh khắc va chạm.
    
    Tại Hoa Kỳ, Cục Quản lý An toàn Giao thông Đường bộ Quốc gia (NHTSA) ước tính rằng lái xe buồn ngủ là nguyên nhân trực tiếp của hơn 100.000 vụ tai nạn mỗi năm, dẫn đến khoảng 1.550 ca tử vong, 71.000 người bị thương, và thiệt hại kinh tế lên tới 12,5 tỷ USD. Các tổ chức khác như CDC cho rằng con số thực tế còn tồi tệ hơn, ước tính có tới 5.000 đến 6.000 vụ tai nạn chết người mỗi năm có yếu tố mệt mỏi. Từ góc độ tổn thất vĩ mô, các vụ va chạm do lái xe buồn ngủ chiếm khoảng 13\% tổng chi phí xã hội khổng lồ trị giá 836 tỷ USD do tai nạn giao thông gây ra tại Mỹ. Tại các quốc gia đang phát triển, con số này cũng ở mức báo động. Các nghiên cứu tại Iran cho thấy mệt mỏi chiếm từ 20\% đến 40\% nguyên nhân tai nạn. Tại Việt Nam, dữ liệu từ Ủy ban An toàn giao thông Quốc gia cho thấy các vụ tai nạn liên quan đến giấc ngủ chiếm tới 30\% tổng số các vụ tai nạn giao thông trong một năm. Một nghiên cứu khác ở vùng nông thôn Ấn Độ cũng chỉ ra rằng 37,6\% tài xế gặp nạn đã không có chế độ nghỉ ngơi đầy đủ trước khi lên xe.

    \begin{table}[htbp]
    \centering
    \caption{Chỉ tiêu thống kê tai nạn do mệt mỏi và buồn ngủ (Dữ liệu toàn cầu \& khu vực)}
    \begin{tabular}{|p{6cm}|p{4cm}|p{4.5cm}|}
    \hline
    \textbf{Chỉ Tiêu Thống Kê} & \textbf{Số Liệu Ước Tính} & \textbf{Ghi chú/Nguồn Trích Dẫn} \\
    \hline
    Tỷ lệ tai nạn toàn cầu do buồn ngủ & 3\% - >30\% & Tùy thuộc khu vực \\
    \hline
    Tỷ lệ tai nạn tại Việt Nam do ngủ gật & 30\% & Ủy ban ATGT Quốc gia \\
    \hline
    Số vụ tai nạn hàng năm tại Mỹ do buồn ngủ & > 100.000 vụ & NHTSA \\
    \hline
    Số người chết hàng năm tại Mỹ do buồn ngủ & 1.550 - 6.000 ca & NHTSA, CDC \\
    \hline
    Thiệt hại kinh tế tại Mỹ do lái xe buồn ngủ & 12,5 tỷ USD & Ước tính hàng năm \\
    \hline
    Tỷ lệ chi phí xã hội tại Mỹ & 13\% tổng chi phí tai nạn (836 tỷ USD) & Tổn thất vĩ mô \\
    \hline
    \end{tabular}
    \end{table}

    \item \textbf{Trạng thái ức chế tâm lý và cảm xúc tiêu cực (Anger/Fear):}\\
    Song song với trạng thái mệt mỏi, các cảm xúc cực đoan như tức giận hoặc hoảng sợ là nguồn cơn của hành vi lái xe hung hăng. Sự căng thẳng khi gặp tắc đường, áp lực thời gian giao hàng, hoặc hành vi ứng xử thiếu văn hóa của các phương tiện xung quanh dễ kích động trạng thái tức giận khi tham gia giao thông (Road Rage). Một người bình thường khi rơi vào trạng thái tâm lý nóng nảy sẽ trải qua sự biến đổi tính cách nhất thời, làm giảm khả năng nhận thức các mối nguy hiểm xung quanh.
    
    Tổ chức AAA Foundation for Traffic Safety báo cáo rằng hành vi lái xe hung hăng chịu trách nhiệm cho hơn 50\% tổng số ca tử vong do tai nạn giao thông. Tốc độ cao là hệ quả trực tiếp phổ biến nhất của trạng thái này, đóng vai trò là yếu tố góp phần trong gần một phần ba số vụ tai nạn xe hơi chết người. Các hành vi điển hình của Road Rage bao gồm: bám đuôi (tailgating) ở khoảng cách không an toàn, lạng lách, chuyển làn liên tục, vượt đèn đỏ, và la mắng phương tiện khác. Cảm xúc chi phối trực tiếp đến hành vi, khả năng phán đoán và thời gian phản ứng của tài xế. Do đó, việc không kiểm soát được cảm xúc tiêu cực không chỉ trực tiếp đe dọa an toàn sinh mạng của bản thân người lái mà còn biến phương tiện thành một cỗ máy gây sát thương vô định đối với cộng đồng.
\end{enumerate}

\subsection{Tầm quan trọng của việc ứng dụng Trí tuệ nhân tạo (AI) trong giám sát lái xe}
Sự cấp thiết phải giải quyết những thách thức do yếu tố con người gây ra đã biến Hệ thống giám sát người lái (Driver Monitoring System - DMS) từ một tính năng thử nghiệm thành một trang bị an toàn mang tính bắt buộc. Trước đây, các giải pháp giám sát tài xế thường dựa vào các phương pháp tiếp cận thụ động hoặc can thiệp vật lý. Điển hình như các cảm biến đo nồng độ cồn qua hơi thở, thiết bị đeo thông minh đo nhịp tim, hoặc các công nghệ Cảnh báo lệch làn (Lane Departure Warning) và đánh giá biến thiên góc vô lăng thuộc nhóm Hỗ trợ người lái tiên tiến (ADAS). Mặc dù đơn giản trong việc triển khai, các hệ thống thụ động này sở hữu nhiều lỗ hổng chí mạng: độ trễ cảnh báo lớn, gây khó chịu cho người lái (khi phải đeo thiết bị liên tục), dễ sinh ra cảnh báo giả (false positives), và quan trọng nhất là không thể đánh giá trực tiếp trạng thái tập trung của tài xế ngay tại thời điểm trước khi chiếc xe có dấu hiệu mất kiểm soát cơ học.

Với sự bứt phá của Trí tuệ nhân tạo (AI) và Thị giác máy tính (Computer Vision), việc ứng dụng các mô hình học sâu (Deep Learning) để phân tích hình ảnh trực tiếp từ camera đặt trong cabin xe (In-cabin Monitoring) đã mở ra một hướng đi mới mang tính cách mạng. Hệ thống DMS hiện đại được vận hành bằng AI có khả năng giải quyết triệt để các hạn chế của công nghệ cũ thông qua ba ưu điểm cốt lõi:
\begin{itemize}
    \item \textbf{Theo dõi không tiếp xúc (Non-intrusive):} Hệ thống chỉ sử dụng một camera RGB hoặc camera hồng ngoại (IR/NIR) phổ thông đặt trên táp-lô, trụ A hoặc gương chiếu hậu để liên tục ghi nhận khuôn mặt tài xế. Môi trường trong khoang lái được số hóa mà không gây bất kỳ cản trở vật lý hay cảm giác khó chịu nào cho người điều khiển phương tiện.
    \item \textbf{Nhận diện đa chiều thời gian thực (Real-time Multi-dimensional Recognition):} Nhờ sức mạnh của các mạng nơ-ron tích chập (CNN) và thuật toán ước lượng mốc hình học khuôn mặt (Face Landmark Detection), hệ thống có thể đồng thời trích xuất các đặc trưng vi mô với độ trễ cực thấp (chỉ vài mili-giây). Nó có khả năng đánh giá mức độ khép mí mắt (PERCLOS), tần suất chớp mắt, biên độ ngáp, hướng nhìn của đồng tử, góc gục đầu, và phân loại sâu các biểu cảm cảm xúc (tức giận, sợ hãi, bình thường) ngay tại thời điểm diễn ra hành vi.
    \item \textbf{Cảnh báo chủ động và Quản trị tập trung:} Hệ thống tự động phân tích và đưa ra tín hiệu beep cảnh báo bằng âm thanh/hình ảnh trực tiếp trong cabin nhằm đánh thức phản xạ của tài xế. Đồng thời, thông qua công nghệ viễn thông (Telematics), dữ liệu về các vi phạm trạng thái được đồng bộ hóa về trung tâm điều hành máy chủ từ xa. Khả năng này giúp các doanh nghiệp vận tải quản lý rủi ro một cách chủ động, can thiệp vào lộ trình trước khi tai nạn xảy ra và tối ưu hóa hệ thống bảo hiểm thương mại.
\end{itemize}

\textbf{Động lực từ Khung pháp lý quốc tế và Thị trường:} \\
Tầm quan trọng của AI trong giám sát khoang lái không chỉ dừng lại ở lý thuyết mà đã được luật hóa ở mức độ toàn cầu. Ủy ban Châu Âu (EU) đã ban hành Quy định An toàn Chung (General Safety Regulation - GSR), bắt buộc tất cả các loại phương tiện mới bán ra tại Châu Âu từ tháng 7 năm 2024 phải trang bị hệ thống Cảnh báo Chú ý và Buồn ngủ (DDAW). Tiến xa hơn, từ năm 2026, quy định Cảnh báo Xao nhãng Nâng cao (ADDW) sẽ bắt buộc sử dụng giám sát chủ động (camera quét mắt) thay vì cảm biến vô lăng thụ động. Tương tự, tổ chức đánh giá an toàn xe hơi độc lập Euro NCAP đã vạch ra lộ trình "Vision 2030", trong đó kể từ năm 2026, các hệ thống giám sát người lái (DMS) là yêu cầu tiên quyết để một chiếc xe đạt được điểm an toàn tối đa (5 sao) trong hạng mục Giám sát trạng thái người ngồi (OSM). Các hệ thống này bắt buộc phải giao tiếp được với hệ thống Phanh khẩn cấp tự động (AEB) để thay đổi mức độ can thiệp khi phát hiện tài xế đang xao nhãng.

\begin{table}[htbp]
\centering
\caption{Khung pháp lý và tiêu chuẩn quốc tế về hệ thống giám sát người lái}
\begin{tabular}{|p{3cm}|p{6cm}|p{5.5cm}|}
\hline
\textbf{Khung Pháp Lý / Tiêu Chuẩn} & \textbf{Quy Định Cốt Lõi} & \textbf{Lộ Trình Áp Dụng} \\
\hline
EU GSR - DDAW & Bắt buộc trang bị hệ thống cảnh báo buồn ngủ và chú ý trên mọi xe thương mại và xe con mới. & 07/07/2022 (Xe mẫu mới) - 07/07/2024 (Toàn bộ xe đăng ký mới) \\
\hline
EU GSR - ADDW & Bắt buộc giám sát chủ động bằng camera phân tích tầm nhìn và mắt thay cho cảm biến thụ động. & Bắt đầu từ 2026 \\
\hline
Euro NCAP (OSM) & Tích hợp hệ thống DMS đánh giá xao nhãng liên kết trực tiếp với ADAS/AEB để duy trì đánh giá an toàn tối đa. & Lộ trình 2023 - 2026 ("Vision 2030") \\
\hline
\end{tabular}
\end{table}

Sức ép pháp lý này đã kích hoạt một chu kỳ tăng trưởng bùng nổ cho thị trường DMS. Các báo cáo phân tích tài chính chỉ ra quy mô thị trường DMS toàn cầu ước tính đạt khoảng 3,5 đến 5,2 tỷ USD vào năm 2025, và được dự phóng sẽ bứt phá lên mức 9,5 đến 12,4 tỷ USD vào giai đoạn 2033 - 2036, với tốc độ tăng trưởng kép hàng năm (CAGR) dao động từ 8,1\% đến 13,3\%. Đặc biệt, các hệ thống dựa trên camera trực quan dự kiến sẽ thống trị thị phần (chiếm trên 43\% đến 57\%) do độ chính xác cao trong việc theo dõi chuyển động mắt và đầu. Việc tích hợp DMS còn mang lại lợi ích tài chính rõ rệt cho các nhà quản lý đội xe; ví dụ, một đội xe 85 chiếc tại Houston (Mỹ) đã tiết kiệm tới 378.000 USD chi phí bảo hiểm thường niên sau khi triển khai hệ thống telematics giám sát hành vi tài xế.

Từ nền tảng lý thuyết sinh học, động lực thị trường, cho đến sự ép buộc của khung pháp lý quốc tế, việc thực hiện đề tài \textbf{"Xây dựng hệ thống nhận dạng biểu cảm khuôn mặt tài xế lái xe, phát hiện cảm xúc tiêu cực, mệt mỏi và hỗ trợ cảnh báo"} là vô cùng cấp thiết. Đề tài không chỉ mang tính thực tiễn cao mà còn là mảnh ghép hoàn hảo để giải quyết triệt để rủi ro từ yếu tố con người, phù hợp với xu hướng chuyển đổi số không thể đảo ngược trong ngành Giao thông thông minh (ITS).

\section{Mục tiêu nghiên cứu và phạm vi nghiên cứu}
\label{section:1.2}
Để chuyển hóa các yêu cầu lý thuyết thành một hệ thống phần mềm có khả năng vận hành thực tế trong điều kiện công nghiệp, giới hạn của đề tài cần được xác định một cách khoa học. Môi trường trong khoang lái ô tô áp đặt các giới hạn cực kỳ nghiêm ngặt về nhiệt độ vận hành, không gian vật lý, năng lượng tiêu thụ, và đặc biệt là năng lực tính toán của hệ thống nhúng.

\subsection{Mục tiêu}
Mục tiêu tối thượng của nghiên cứu là thiết kế và xây dựng một giải pháp lai (Hybrid System) kết hợp tối ưu giữa hai trường phái tính toán của Thị giác máy tính: (1) Phân tích các đặc trưng hình học khuôn mặt dựa trên mốc không gian 3D để nhận diện trạng thái buồn ngủ, mất tập trung; và (2) Ứng dụng mô hình mạng học sâu phân loại ảnh tích chập (CNN) để nhận diện các biểu cảm cảm xúc (tức giận, sợ hãi).

Hệ thống hướng tới việc đạt được độ tin cậy cao về mặt nhận diện, đồng thời phải thỏa mãn tiêu chí hoạt động mượt mà, thời gian thực (real-time) trên các thiết bị máy tính biên (Edge Computing Devices) có tài nguyên phần cứng giới hạn (CPU tiêu chuẩn, không phụ thuộc vào các vi mạch xử lý đồ họa GPU đắt tiền). Sự kết hợp này nhằm loại bỏ các cảnh báo giả, tối ưu hóa băng thông truyền tải và xây dựng một quy trình quản trị rủi ro hoàn thiện từ Thiết bị đầu cuối (Client) đến Máy chủ đám mây (Cloud Server).

\subsection{Đối tượng nghiên cứu}
Hệ sinh thái công nghệ cấu thành nên giải pháp DMS đòi hỏi sự tập trung vào các đối tượng nghiên cứu cụ thể, bao gồm kiến trúc mạng nơ-ron, thư viện trích xuất đặc trưng, và nguồn dữ liệu huấn luyện.
\begin{enumerate}
    \item \textbf{Các mô hình học sâu phân loại ảnh đa lớp:}\\
    Nghiên cứu đặt trọng tâm vào việc đánh giá và tinh chỉnh kiến trúc MobileNetV3 (phiên bản Large). So với các mạng học sâu khổng lồ như VGG19 (tuy đạt độ chính xác lên đến 99,7\% trên tập dữ liệu KMU-FED nhưng lại đòi hỏi hàng tỷ phép tính dấu phẩy động FLOPs), MobileNetV3 mang tính đột phá cho môi trường thiết bị di động. Nó tận dụng các khối residual đảo ngược (inverted residual blocks), tích hợp thuật toán tìm kiếm cấu trúc mạng (Network Architecture Search) và cơ chế Attention nhẹ (Squeeze-and-Excitation). Kiến trúc này cho phép suy luận tốc độ cao, giữ được lượng tham số nhỏ mà vẫn đảm bảo độ chính xác vượt trội khi phân loại các lớp cảm xúc phức tạp. Đi kèm với đó, mô hình BlazeFace được nghiên cứu để giải bài toán định vị và bám sát khuôn mặt (Face Detection) trong các điều kiện thiếu sáng hoặc bị che khuất.
    
    \item \textbf{Các giải pháp trích xuất mốc hình học không gian:}\\
    Đối tượng nghiên cứu trọng yếu cho bài toán buồn ngủ là thư viện MediaPipe FaceMesh. Giải pháp này cung cấp khả năng trích xuất ma trận tọa độ không gian 3D của 468 mốc khuôn mặt (landmarks) theo thời gian thực. Từ ma trận tọa độ này, hệ thống có thể tính toán chính xác biên độ mở của mí mắt và khẩu hình miệng thông qua các phương trình hình học Euclid mà không cần qua các lớp xử lý CNN nặng nề.
    
    \item \textbf{Tập dữ liệu huấn luyện (Dataset Ecosystem):}\\
    Sức mạnh của một thuật toán AI bị chi phối trực tiếp bởi chất lượng và tính đại diện của dữ liệu huấn luyện. Không gian cabin là một môi trường quang học khắc nghiệt: ánh sáng thay đổi liên tục (đi vào hầm chui, ngược nắng, ban đêm), các che khuất vật lý (kính râm, khẩu trang), và sự đa dạng chủng tộc của tài xế. Do đó, nghiên cứu khảo sát và tái cấu trúc một hệ sinh thái dữ liệu bao gồm các bộ dữ liệu khuôn mặt tổng quát và các bộ dữ liệu chuyên biệt cho môi trường khoang lái ô tô:
    \begin{itemize}
        \item \textbf{Các tập dữ liệu về trạng thái buồn ngủ/xao nhãng:}\\
        NTHU-DDD (National Tsing Hua University Drowsy Driver Detection): Chứa 9,5 giờ video phân mảnh từ 18 đối tượng, đánh nhãn các trạng thái buồn ngủ, ngáp và mất tập trung trong bối cảnh phòng thí nghiệm giả lập.\\
        UTA-RLDD (Real-Life Drowsiness Dataset): Ghi lại quá trình chuyển giao nhận thức tự nhiên của tài xế từ trạng thái tỉnh táo sang mệt mỏi nhẹ và cuối cùng là ngủ gật cực độ. Tập dữ liệu được phân chia theo các cửa sổ thời gian 30 giây để phục vụ phân tích chuỗi trạng thái liên tục.
        
        \item \textbf{Các tập dữ liệu về cảm xúc và đặc điểm chiếu sáng:}\\
        MLI-DER (Multiple Light Intensities Driver Emotion Recognition): Một cơ sở dữ liệu đột phá chuyên giải quyết bài toán nhiễu sáng (Illumination invariance). Tập dữ liệu cung cấp các hình ảnh khuôn mặt tài xế dưới nhiều cường độ sáng phức tạp, giúp mô hình bóc tách đặc trưng cảm xúc khỏi yếu tố ngoại cảnh.\\
        KMU-FED (Keimyung University Facial Expression of Drivers): Cung cấp các chuỗi hình ảnh thực tế thu bằng camera cận hồng ngoại (NIR) gắn trên vô lăng/táp-lô từ 12 đối tượng lái xe thực tế. Bộ dữ liệu này có ý nghĩa quan trọng trong việc huấn luyện mô hình vượt qua các rào cản về bóng đổ hình học và sự che khuất do kính phân cực hay góc quay chéo.
    \end{itemize}
    Ngoài ra, các tập dữ liệu biểu cảm đại chúng như RAF-DB, AffectNet, FER-2013, KDEF và SFEW được nghiên cứu sử dụng làm trọng số cơ sở (baseline weights) để áp dụng kỹ thuật Học chuyển giao (Transfer Learning) trước khi tinh chỉnh riêng cho không gian cabin.
\end{enumerate}

\subsection{Phạm vi nghiên cứu: Luồng dữ liệu video trong cabin xe}
Để đảm bảo tính khả thi và tập trung sâu vào giải quyết cốt lõi bài toán, phạm vi nghiên cứu được khoanh vùng ở các khía cạnh không gian, chức năng và phần cứng như sau:
\begin{itemize}
    \item \textbf{Vị trí thu nhận tín hiệu:} Hệ thống chỉ xử lý luồng dữ liệu video trực tiếp thu nhận từ góc nhìn trực diện phía trước tài xế. Môi trường giả lập tương đương với việc bố trí một camera hành trình hoặc webcam đặt tại khu vực gương chiếu hậu trung tâm hoặc trên bảng điều khiển táp-lô (dashboard), mô phỏng góc nhìn quang học bao quát phần đầu và vai của tài xế.
    
    \item \textbf{Phạm vi nhận dạng và phân loại:}
    \begin{itemize}
        \item \textbf{Nhóm 1:} Nhận dạng trạng thái thể chất suy giảm bao gồm mệt mỏi, buồn ngủ (Fatigue/Drowsiness) và hành vi xao nhãng hướng nhìn (Distracted) dẫn đến mất tập trung trên đường.
        \item \textbf{Nhóm 2:} Phân loại 5 lớp cảm xúc cốt lõi có tác động trực tiếp đến động năng lái xe, bao gồm: trạng thái Bình thường (Neutral - baseline), Tức giận (Anger - nguyên nhân Road Rage), Sợ hãi (Fear - khủng hoảng tâm lý), Vui vẻ (Happiness) và Buồn bã (Sadness).
    \end{itemize}
    
    \item \textbf{Giới hạn nền tảng phần cứng:} Việc thiết kế mô hình bắt buộc phải tối ưu hóa để thực thi quá trình suy luận (Inference) thời gian thực trên các bộ vi xử lý trung tâm (CPU) thông thường. Hệ thống không yêu cầu thiết bị xử lý phải có card đồ họa chuyên dụng (GPU cao cấp) ở phía Client. Giới hạn này đảm bảo tính khả thi cao nhất khi triển khai hàng loạt dự án ra thực tế trên các hộp đen viễn thông (Telematics Box) hoặc thiết bị máy tính biên giá rẻ như Raspberry Pi gắn trên hàng triệu phương tiện thương mại.
\end{itemize}

\section{Nhiệm vụ đồ án}
\label{section:1.3}
Để giải quyết triệt để bài toán an toàn giao thông nêu trên và hoàn thành các mục tiêu khoa học, đồ án xác định rõ 5 nhiệm vụ thực thi cụ thể với độ chi tiết cao:
\begin{enumerate}
    \item \textbf{Nghiên cứu lý thuyết và cơ sở toán học sinh trắc:}\\
    Nhiệm vụ đầu tiên là tìm hiểu sâu về nền tảng cơ sinh học và phương pháp số hóa chuyển động khuôn mặt. Tập trung nghiên cứu công thức toán học tính toán các chỉ số vi mô như EAR (Eye Aspect Ratio) để đánh giá biên độ mở của nhãn cầu, và MAR (Mouth Aspect Ratio) để đo lường động học của việc ngáp. Về mặt lý thuyết, EAR được tính toán bằng tỷ lệ khoảng cách dọc và ngang giữa các mốc điểm quanh mắt. Mức EAR bình thường duy trì trên 0.20, và sụt giảm về gần 0 khi nhắm mắt. Hệ thống cần phát hiện các khoảnh khắc EAR rơi xuống dưới ngưỡng cảnh báo (ví dụ 0.15 hoặc 0.20). Tương tự, MAR > 0.60 thường được sử dụng làm mốc xác định một hành vi ngáp. Hơn thế nữa, nghiên cứu mở rộng sang các chỉ số vĩ mô như góc nghiêng trục đầu (Pitch) và tiêu chuẩn PERCLOS (Percentage of Eyelid Closure). PERCLOS đánh giá tỷ lệ phần trăm thời gian mí mắt nhắm vượt quá 80\% trong một chu kỳ thời gian nhất định (ví dụ 1 phút), được y khoa công nhận là chỉ báo độ tin cậy cao nhất cho mức độ kiệt sức. Đồng thời, đi sâu vào lý thuyết kiến trúc mạng CNN tối ưu dành cho thiết bị di động MobileNetV3 để chứng minh sự vượt trội về thông lượng xử lý.
    
    \item \textbf{Xây dựng và chuẩn hóa bộ dữ liệu tích hợp:}\\
    Thực hiện quá trình thu thập, tiền xử lý và đồng nhất hóa cấu trúc dữ liệu từ các bộ dữ liệu khuôn mặt và trạng thái lái xe rời rạc trên thế giới (như RAF-DB, NTHU DDD, UTA-RLDD, MLI-DER, KMU-FED). Nhiệm vụ này yêu cầu xây dựng các kịch bản ánh xạ nhãn (label mapping), loại bỏ các dữ liệu nhiễu, cân bằng số lượng mẫu giữa các lớp cảm xúc, và tái cấu trúc chúng thành một tập dữ liệu chuẩn DMS duy nhất phục vụ vòng đời huấn luyện mô hình.
    
    \item \textbf{Huấn luyện mô hình trí tuệ nhân tạo:}\\
    Lập trình quy trình tiền xử lý hình ảnh nâng cao nhằm chống lại các biến động môi trường. Kỹ thuật cân bằng lược đồ mức xám cục bộ CLAHE (Contrast Limited Adaptive Histogram Equalization) được áp dụng để khôi phục chi tiết vùng tối. Kết hợp với kỹ thuật cắt cúp khuôn mặt tự động (Face Cropping) để tập trung ma trận ảnh vào vùng chứa đặc trưng. Tiếp theo, thiết lập môi trường bằng framework PyTorch, định nghĩa hàm mất mát (Loss Function), bộ tối ưu hóa (Optimizer) và thực thi quá trình huấn luyện mô hình MobileNetV3. Các siêu tham số (hyperparameters) được tinh chỉnh lặp lại để đạt điểm cân bằng giữa độ chính xác dự đoán (Accuracy) và tránh hiện tượng học vẹt (Overfitting).
    
    \item \textbf{Xây dựng chương trình phần mềm thực nghiệm Edge-to-Cloud:}\\
    Nhiệm vụ trọng tâm là hiện thực hóa thuật toán thành một hệ thống phần mềm hoàn chỉnh, được chia làm hai phân hệ tương tác hai chiều:
    \begin{itemize}
        \item \textbf{Phía Client (AI Biên):} Lập trình module đọc luồng video từ camera sử dụng kỹ thuật xử lý đa luồng (multi-threading) để tránh nghẽn bộ đệm. Tích hợp engine suy luận AI, kết hợp các bộ lọc tín hiệu thời gian thực nhằm khử nhiễu dữ liệu. Khi phát hiện vi phạm, phần mềm lập tức kích hoạt các module phát âm thanh (beep) cảnh báo cục bộ tới người lái.
        \item \textbf{Phía Server (Trung tâm):} Thiết lập một máy chủ ứng dụng web (Web Server) tối giản và hiệu suất cao bằng framework Flask. Xây dựng các cổng giao tiếp API (RESTful APIs) để lắng nghe và thu nhận thông tin viễn thông. Dữ liệu được cấu trúc và lưu trữ vĩnh viễn vào hệ quản trị cơ sở dữ liệu SQLite. Phát triển cơ chế đồng bộ hóa dữ liệu từ máy chủ đẩy thẳng lên giao diện điều hành (Dashboard) trên trình duyệt web thông qua luồng kết nối liên tục Server-Sent Events (SSE).
    \end{itemize}
    
    \item \textbf{Thử nghiệm thực địa và đánh giá hệ thống:}\\
    Triển khai phần mềm lên các cấu hình phần cứng thử nghiệm để đánh giá toàn diện. Các tiêu chí đánh giá không chỉ dừng ở độ chính xác phân loại của mô hình học sâu (Precision, Recall, F1-Score) trên tập kiểm thử tĩnh, mà còn đo lường hiệu năng kỹ thuật của hệ thống trên luồng camera động thời gian thực. Các thông số như Tốc độ xử lý khung hình trên giây (FPS), mức tiêu hao năng lượng CPU, và độ trễ truyền tin được ghi nhận. Đồng thời, phân tích sâu các kịch bản lỗi ở vùng biên (edge cases) như tài xế xoay mặt quá góc độ cho phép hoặc thay đổi cường độ sáng đột ngột, từ đó đề xuất phương án tinh chỉnh cho các phiên bản sau.
\end{enumerate}

\section{Định hướng giải pháp}
\label{section:1.4}
Hệ thống DMS được định hướng thiết kế theo mô hình kiến trúc lai phân tán Edge-to-Cloud, phản ánh xu thế hiện đại trong thiết kế các hệ thống Internet vạn vật (IoT). Mô hình này giải quyết triệt để vấn đề nghẽn cổ chai mạng bằng cách đẩy năng lực xử lý tính toán ra sát với nguồn sinh dữ liệu (biên), và chỉ sử dụng mạng viễn thông (4G/5G) để truyền tải các tín hiệu siêu dữ liệu (metadata) mang tính quyết định. Hệ thống được cấu thành từ hai phân hệ độc lập nhưng liên kết chặtẽ:
\begin{enumerate}
    \item \textbf{Phân hệ AI Biên (Edge Client) và Thuật toán tối ưu tính toán:}\\
    Phân hệ phía khách (Client) bao gồm một webcam tiêu chuẩn kết nối trực tiếp với thiết bị máy tính nhúng (hoặc laptop giả lập) gắn bên trong xe. Khác với các hệ thống phụ thuộc đám mây, toàn bộ phân hệ này chạy mã nguồn Python biên dịch tại chỗ, sử dụng cơ chế xử lý đa luồng (Multi-threading). Cơ chế này tách biệt quá trình đọc ảnh (I/O bounds) và quá trình xử lý AI (CPU bounds), đảm bảo khung hình camera luôn được bắt giữ liên tục mà không gặp hiện tượng giật lag hay trễ hình.
    
    Để giải quyết bài toán giới hạn tài nguyên CPU, giải pháp ứng dụng thiết kế xử lý bất đối xứng theo chu kỳ, gọi là Bỏ qua khung hình (Skip-frame). Việc trích xuất các mốc 3D và tính toán EAR/MAR bằng MediaPipe vốn cực kỳ nhẹ bén, do đó sẽ được thực thi trên mọi khung hình (toàn thời gian). Tuy nhiên, mạng CNN học sâu MobileNetV3 (phục vụ việc phân loại 5 trạng thái cảm xúc phức tạp) lại tiêu thụ nhiều nhiệt năng tính toán. Do trạng thái cảm xúc của con người hiếm khi biến đổi hoàn toàn trong khoảng thời gian dưới 1 giây, hệ thống được thiết lập để chỉ chạy suy luận CNN ngắt quãng (ví dụ: chỉ chạy 1 lần sau mỗi 5 hoặc 10 khung hình). Kỹ thuật này giữ cho tốc độ FPS tổng thể luôn ở mức cao mà vẫn bảo toàn tính liên tục của việc nhận thức ngữ cảnh hành vi.
    
    Khi thuật toán tại Client phát hiện chỉ số EAR, MAR hoặc cảm xúc vi phạm ngưỡng an toàn đã cấu hình, phần mềm sẽ không lập tức phát tín hiệu đi. Đây là lúc logic cốt lõi tiếp theo tham gia định hướng giải pháp: Bộ lọc Đảo nhiễu (Debounce Logic). Trong thực tế đo lường tín hiệu chuỗi thời gian, các giá trị thường dao động hỗn loạn ở sát vạch ngưỡng. Một cái chớp mắt thông thường hay một cái tặc lưỡi có thể khiến EAR/MAR chạm ngưỡng trong 1-2 khung hình. Nếu không có bộ lọc, hệ thống sẽ kích hoạt vô số cảnh báo giả, gây ra hội chứng "mệt mỏi vì cảnh báo" (alert fatigue) cho tài xế và làm tràn bộ đệm cơ sở dữ liệu. Cơ chế Debounce xử lý vấn đề này bằng cách áp dụng một cửa sổ thời gian chờ (delay window) khoảng từ 800 mili-giây đến 1 giây. Khi phát hiện một lỗi vượt ngưỡng (ví dụ MAR chỉ ra trạng thái "Ngáp"), hệ thống khởi động bộ đếm thời gian. Nếu trạng thái "Ngáp" không được duy trì liên tục và biến mất trước khi hết khoảng thời gian $T_{debounce}$, nó bị phân loại là nhiễu và bị triệt tiêu. Chỉ khi tín hiệu ổn định vượt quá bộ đếm thời gian, sự kiện mới chính thức được ghi nhận. Lúc này, Client sẽ xuất ra tín hiệu âm thanh Beep cảnh báo cục bộ qua hệ thống loa xe, đồng thời đóng gói dữ liệu cấu trúc dạng JSON và thực hiện một lệnh gọi POST lên Server.
    
    \item \textbf{Phân hệ Quản trị Trung tâm (Cloud Server) và Đồng bộ hóa thời gian thực:}\\
    Tại trung tâm điều hành, máy chủ đám mây triển khai nền tảng Web Server tối giản bằng Flask. Nhiệm vụ của máy chủ là lắng nghe các yêu cầu POST liên tục cập nhật trạng thái từ hàng ngàn phương tiện gửi về. Nhờ cơ chế Debounce tại Client đã lọc sạch nhiễu, lượng dữ liệu đổ về Server là các chuỗi cảnh báo đặc và có giá trị cao, giúp giảm thiểu đáng kể xung đột ghi vào hệ quản trị cơ sở dữ liệu quan hệ cục bộ SQLite. SQLite được chọn nhờ tính bảo mật cao, khả năng lưu trữ tệp tin vĩnh viễn không cần cấu hình phức tạp, hoàn toàn đáp ứng được quy mô dữ liệu của các đội xe tầm trung.
    
    Định hướng giải pháp cho giao diện người dùng (Dashboard) là triệt tiêu hoàn toàn độ trễ hiển thị. Thay vì yêu cầu trình duyệt của người quản lý phải liên tục tải lại trang (polling - một giải pháp cổ điển gây quá tải máy chủ), hệ thống thiết lập một luồng Server-Sent Events (SSE). Phương thức luồng một chiều (unidirectional stream) này duy trì một đường ống kết nối HTTP mở vĩnh viễn giữa Server và Client Web. Ngay tại thời điểm cơ sở dữ liệu SQLite ghi nhận một sự kiện tai nạn hoặc biểu cảm tức giận (Anger) từ bất kỳ xe nào, máy chủ lập tức đẩy (push) gói cập nhật đó xuống thẳng trình duyệt. Giao diện frontend sẽ thu nhận gói tin và tái hiện tức thời danh sách phương tiện vi phạm, cập nhật đồ thị Chart.js và tính toán lại điểm số rủi ro của tài xế. Nhờ đó, doanh nghiệp vận tải sở hữu một tháp không lưu viễn thông thực thụ, trao quyền can thiệp vào lộ trình của tài xế trước khi những biểu hiện tâm lý và thể chất kia chuyển hóa thành các tai nạn thảm khốc.
\end{enumerate}

\section{Bố cục đồ án}
\label{section:1.5}
Dựa trên các mục tiêu và định hướng giải pháp đã thiết lập, cấu trúc của báo cáo đồ án tốt nghiệp được chia nhỏ và hệ thống hóa thành 6 chương chuyên sâu. Mỗi chương đóng vai trò giải quyết một khía cạnh riêng biệt từ nền tảng lý thuyết cho đến thực thi kỹ thuật, đảm bảo tính liền mạch của một công trình nghiên cứu khoa học:
\begin{itemize}
    \item \textbf{Chương 1. Giới thiệu đề tài:} Cung cấp bức tranh toàn cảnh về vấn đề nghiên cứu. Trình bày các dữ liệu dịch tễ học và luật pháp để làm nổi bật lý do chọn đề tài và tính cấp thiết. Định hình rõ mục tiêu, đối tượng, phạm vi nghiên cứu (giới hạn ở luồng video cabin, tập trung mô hình MobileNetV3) và phác thảo định hướng kiến trúc giải pháp Edge-to-Cloud.
    \item \textbf{Chương 2. Khảo sát và phân tích yêu cầu hệ thống:} Đi sâu vào phân tích hệ sinh thái công nghệ hiện hữu. Khảo sát chi tiết tính chất của các bộ dữ liệu khuôn mặt dùng để huấn luyện (đặc biệt là NTHU DDD, UTA-RLDD, MLI-DER và KMU-FED). Căn cứ vào đó, tiến hành lập tài liệu đặc tả yêu cầu, phân tách rạch ròi các yêu cầu chức năng (các loại cảnh báo, giao tiếp API) và yêu cầu phi chức năng (FPS tối thiểu, mức tiêu hao tài nguyên, tính ổn định) của toàn bộ hệ thống.
    \item \textbf{Chương 3. Nền tảng lý thuyết và công nghệ sử dụng:} Xây dựng khung lý luận vững chắc cho quá trình lập trình. Giới thiệu cơ sở toán học và phương pháp luận chiết xuất hình học không gian của các chỉ số sinh trắc học EAR, MAR, tiêu chuẩn đo lường mệt mỏi PERCLOS. Diễn giải lý thuyết kiến trúc mạng học sâu MobileNetV3 (các khối bottleneck, chiều sâu layer) và phân tích vai trò của các công nghệ phần mềm phụ trợ (Framework PyTorch, môi trường Flask, cơ sở dữ liệu SQLite, và giao thức kết nối SSE).
    \item \textbf{Chương 4. Phân tích thiết kế mô hình và hệ thống:} Thể hiện tư duy kỹ sư phần mềm bằng việc thiết kế bản vẽ kiến trúc. Trình bày chi tiết luồng xử lý (Pipeline) tổng thể qua các lưu đồ thuật toán (Flowcharts). Đặc tả quá trình thiết kế lược đồ cơ sở dữ liệu (Database Schema) trong SQLite, thiết kế các mô-đun tiền xử lý ảnh (CLAHE, tăng cường dữ liệu), cũng như mô hình hóa toán học cho logic ra quyết định cảnh báo và cơ chế Đảo nhiễu (Debounce filter).
    \item \textbf{Chương 5. Xây dựng chương trình, thử nghiệm và đánh giá:} Trình bày nhật ký quá trình thực thi từ mã nguồn lý thuyết đến sản phẩm thực tế. Mô tả quá trình huấn luyện mô hình học sâu (đồ thị Loss/Accuracy), các giải pháp lập trình xử lý đa luồng (Multi-threading), và kỹ thuật lặp bỏ qua khung hình (Skip-frame) để tối ưu hệ thống thiết bị biên. Cung cấp hệ thống bảng biểu, số liệu thực nghiệm đo đạc đánh giá độ chính xác của mô hình và hiệu năng tổng thể của phần mềm (như chỉ số FPS/tiêu thụ điện năng CPU).
    \item \textbf{Chương 6. Kết luận và hướng phát triển:} Đóng gói lại công trình bằng cách tổng kết các kết quả khoa học đã đạt được, tiến hành đối chiếu nghiêm ngặt với các mục tiêu ban đầu trong phiếu giao nhiệm vụ đồ án. Trình bày minh bạch các điểm hạn chế còn tồn tại (edge cases, giới hạn điều kiện thiếu sáng cực độ) và mở ra các định hướng nâng cấp đầy hứa hẹn trong tương lai (như tích hợp sâu module DMS vào hệ thống điều khiển phanh tự động của xe tự hành).
\end{itemize}

\end{document}
